\heading{实验物理中的统计方法-25春-期末}
\noteinfo{内容精修版；已统一“拒绝/不拒绝”表述，并修正题干中缺字与参数条件。}

\begin{question}
在假设检验中，零假设为 $H_0$。记
\[
P_1=P(\text{拒绝 }H_0\mid H_0\text{ 为真}),
\qquad
P_2=P(\text{不拒绝 }H_0\mid H_0\text{ 为假}),
\]
则 $P_1$ 被称为弃真错误，$P_2$ 被称为取伪错误。
\begin{solution}
$P_1$ 是在 $H_0$ 为真时拒绝 $H_0$ 的概率，即第一类错误，也称弃真错误；
$P_2$ 是在 $H_0$ 为假时不拒绝 $H_0$ 的概率，即第二类错误，也称取伪错误。故命题正确。
\anstext{正确}
\end{solution}
\end{question}

\begin{question}
设总体 $X$ 的均值和方差分别为 $\mu$ 和 $\sigma^2$，$X_1,X_2,\cdots,X_n$ 是来自该总体的一个简单随机样本。仅由这些条件，能否断定 $\bar X=\dfrac1n\sum_{i=1}^n X_i$ 一定是 $\mu$ 的有效估计量？
\begin{solution}
样本均值
\[
\bar X=\frac1n\sum_{i=1}^nX_i
\]
确实满足
\[
E[\bar X]=\mu,
\qquad
\V(\bar X)=\frac{\sigma^2}{n},
\]
因此它是 $\mu$ 的无偏估计量和矩估计量。但是“有效估计量”通常指在无偏估计类中方差达到 Cram\'er--Rao 下界或具有最小方差的估计量。题干只给出总体均值和方差，没有给出总体分布或 Fisher 信息，不能推出 $\bar X$ 一定有效。故不能仅由题设断定其一定有效。
\anstext{不能；原命题严格说错误}
\end{solution}
\end{question}

\begin{question}
假设 $X_1,\cdots,X_n$ 相互独立且均服从标准正态分布，则 $\sum_{i=1}^nX_i^2$ 服从自由度为 $n$ 的卡方分布。
\begin{solution}
卡方分布的定义即为 $n$ 个相互独立的标准正态随机变量平方和的分布。因此
\[
\sum_{i=1}^nX_i^2\sim \chi^2(n).
\]
故命题正确。
\anstext{正确}
\end{solution}
\end{question}

\begin{question}
设 $X\sim\chi^2(2n)$，则当 $n$ 很大时，$\dfrac{X-2n}{2\sqrt n}$ 近似服从标准正态分布。
\begin{solution}
若 $X\sim\chi^2(2n)$，则
\[
E[X]=2n,
\qquad
\V(X)=2(2n)=4n.
\]
卡方分布在自由度较大时渐近正态，因此
\[
\frac{X-2n}{2\sqrt n}\xrightarrow[]{d}N(0,1).
\]
故命题正确。
\anstext{正确}
\end{solution}
\end{question}

\begin{question}
假设 $X_1,X_2,\cdots,X_n$ 是总体 $X$ 的一个简单随机样本，则估计量
\[
\frac{1}{n-1}\sum_{i=1}^n(X_i-\bar X)^2
\]
是方差 $\sigma^2$ 的无偏估计。
\begin{solution}
由样本方差的基本性质，
\[
E\left[\sum_{i=1}^n(X_i-\bar X)^2\right]=(n-1)\sigma^2.
\]
因此
\[
E\left[\frac{1}{n-1}\sum_{i=1}^n(X_i-\bar X)^2\right]=\sigma^2.
\]
故命题正确。
\anstext{正确}
\end{solution}
\end{question}

\begin{question}
设 $X\sim N(\mu,\sigma^2)$，其中 $\sigma^2$ 已知。若样本容量 $n$ 和置信水平 $1-\alpha$ 均保持不变，则对于不同的样本观测值，参数 $\mu$ 的等尾双侧置信区间的长度保持不变。
\begin{solution}
当 $\sigma^2$ 已知时，$\mu$ 的等尾双侧置信区间为
\[
\left[\bar X-z_{1-\alpha/2}\frac{\sigma}{\sqrt n},\
\bar X+z_{1-\alpha/2}\frac{\sigma}{\sqrt n}\right].
\]
其长度为
\[
2z_{1-\alpha/2}\frac{\sigma}{\sqrt n},
\]
只依赖于 $n,\alpha,\sigma$，不依赖于具体观测到的 $\bar X$。故命题正确。
\anstext{正确}
\end{solution}
\end{question}

\begin{question}
设总体 $X$ 服从参数为 $(n,p)$ 的二项分布 $(0<p<1,n\text{ 为正整数})$，$(X_1,X_2,\cdots,X_N)$ 为来自该总体的简单随机样本，$\bar X$ 为样本均值，则
\[
E\left[\sum_{i=1}^N(X_i-\bar X)^2\right]=n(N-1)p(1-p).
\]
\begin{solution}
对单个观测量，
\[
\V(X_i)=np(1-p).
\]
又因为
\[
E\left[\sum_{i=1}^N(X_i-\bar X)^2\right]=(N-1)\V(X_i),
\]
所以
\[
E\left[\sum_{i=1}^N(X_i-\bar X)^2\right]
=(N-1)np(1-p).
\]
故命题正确。
\anstext{正确}
\end{solution}
\end{question}

\begin{question}
设 $(X_1,X_2,\cdots,X_n)$ 是来自正态总体 $N(\mu,\sigma^2)$ 的简单随机样本，则 $\mu$ 的矩估计量为
\[
\hat\mu=\frac1n\sum_{i=1}^nX_i.
\]
\begin{solution}
矩估计法用样本一阶矩代替总体一阶矩。由于
\[
E[X]=\mu,
\qquad
\frac1n\sum_{i=1}^nX_i=\bar X,
\]
所以
\[
\hat\mu_{\mathrm{MOM}}=\bar X.
\]
故命题正确。
\anstext{正确}
\end{solution}
\end{question}

\begin{question}
设总体 $X$ 的均值为 $\mu$，方差有限。$\theta_1$ 和 $\theta_2$ 是参数 $\mu$ 的两个无偏估计量，其方差分别为 $\sigma_1^2$ 和 $\sigma_2^2$，相关系数为 $\rho$。已知
\[
\theta=c_1\theta_1+c_2\theta_2
\]
是在所有满足 $c_1+c_2=1$ 的线性无偏组合中方差最小的估计量。若最优解满足 $c_1>0,c_2>0$，则
\[
c_1=\blank[5cm],\qquad c_2=\blank[5cm].
\]
\begin{solution}
由 $c_1+c_2=1$，令 $c_2=1-c_1$。记
\[
\cov(\theta_1,\theta_2)=\rho\sigma_1\sigma_2.
\]
则
\[
\V(\theta)
=c_1^2\sigma_1^2+(1-c_1)^2\sigma_2^2
+2c_1(1-c_1)\rho\sigma_1\sigma_2.
\]
对 $c_1$ 求导并令其为零，得
\[
2c_1\sigma_1^2-2(1-c_1)\sigma_2^2
+2(1-2c_1)\rho\sigma_1\sigma_2=0.
\]
整理得
\[
c_1(\sigma_1^2+\sigma_2^2-2\rho\sigma_1\sigma_2)
=\sigma_2^2-\rho\sigma_1\sigma_2.
\]
故
\[
c_1=\frac{\sigma_2^2-\rho\sigma_1\sigma_2}
{\sigma_1^2+\sigma_2^2-2\rho\sigma_1\sigma_2},
\qquad
c_2=\frac{\sigma_1^2-\rho\sigma_1\sigma_2}
{\sigma_1^2+\sigma_2^2-2\rho\sigma_1\sigma_2}.
\]
等价地，
\ansmath{c_1=\frac{\sigma_2(\sigma_2-\rho\sigma_1)}
{\sigma_1^2-2\rho\sigma_1\sigma_2+\sigma_2^2},
\qquad
c_2=\frac{\sigma_1(\sigma_1-\rho\sigma_2)}
{\sigma_1^2-2\rho\sigma_1\sigma_2+\sigma_2^2}}
\end{solution}
\end{question}

\begin{question}
设 $(X_1,X_2,\cdots,X_n)$ 是来自正态总体 $N(\mu,\sigma^2)$ 的简单随机样本，参数 $\mu,\sigma^2$ 均未知，样本均值 $\bar X=2025$。若已知参数 $\mu$ 的置信水平为 $95\%$ 的等尾双侧置信区间的上侧界限为 $2050$，则 $\mu$ 的置信水平为 $95\%$ 的等尾双侧置信区间是 \blank[4cm]。
\begin{solution}
等尾双侧区间关于样本均值 $\bar X$ 对称。已知中心为 $2025$，上侧界限为 $2050$，则半长为
\[
2050-2025=25.
\]
因此下侧界限为
\[
2025-25=2000.
\]
故置信区间为
\ansmath{[2000,2050]}
\end{solution}
\end{question}

\begin{question}
轴子是暗物质的候选者之一。若某类轴子存在，探测器中会出现相应的候选事件。假设某实验预期的本底数为 $3$ 个，实际观测到疑似信号 $5$ 个。试计算该观测结果的显著性水平。已知泊松分布
\[
f(n;\nu)=\frac{\nu^n e^{-\nu}}{n!}.
\]
\begin{solution}
零假设取为只有本底，即
\[
H_0:\nu=3.
\]
观测到 $5$ 个或更多事件越偏离本底假设，因此显著性水平取右尾概率
\[
\alpha=P(N\ge 5\mid \nu=3).
\]
于是
\[
\alpha=1-P(N\le 4\mid \nu=3)
=1-\sum_{n=0}^4\frac{3^ne^{-3}}{n!}.
\]
展开为
\[
\alpha=1-e^{-3}\left(1+3+\frac{3^2}{2!}+\frac{3^3}{3!}+\frac{3^4}{4!}\right).
\]
数值上
\[
\alpha\approx 0.1847.
\]
因此该观测结果的显著性水平约为 $0.185$。
\ansmath{\alpha=1-\sum_{n=0}^4 f(n;3)\approx 0.185}
\end{solution}
\end{question}

\begin{question}
假设实验上观测某稀有事件，在给定的观测时间内未观测到事件数 $n_{\mathrm{obs}}=0$，且预期本底数 $\mu_{\mathrm{bkg}}$ 可以忽略。求预期信号数 $\mu_{\mathrm{sig}}$ 在 $95\%$ 置信水平下的上限。
\begin{solution}
本底可忽略时，观测数服从
\[
N\sim \Pois(\mu_{\mathrm{sig}}).
\]
若 $95\%$ 上限为 $\mu_{\mathrm{up}}$，则
\[
P(N=0\mid \mu_{\mathrm{up}})=1-0.95=0.05.
\]
由泊松分布
\[
P(N=0\mid \mu_{\mathrm{up}})=e^{-\mu_{\mathrm{up}}},
\]
故
\[
e^{-\mu_{\mathrm{up}}}=0.05.
\]
因此
\[
\mu_{\mathrm{up}}=-\ln 0.05=2.996\approx 3.
\]
\ansmath{\mu_{\mathrm{sig}}^{\mathrm{up}}=-\ln(1-0.95)=2.996\approx 3}
\end{solution}
\end{question}

\begin{question}
假设有一对变量 $(x,y)$，例如 $x$ 可以表示样品的温度，$y$ 可以表示样品的热容。假定变量 $y$ 服从正态分布 $N(\mu,\sigma^2)$，方差 $\sigma^2$ 保持不变，而均值
\[
\mu=\alpha+\beta(x-\bar x),
\]
其中 $\bar x$ 为样本均值，且 $x$ 不一定是随机变量。即在 $n$ 次相互独立的测量中，在指定的 $x=x_i$ 处观测得到 $y_i$，$i=1,2,\cdots,n$。求参数 $\alpha,\beta$ 的最大似然估计量。
\begin{solution}
对给定的 $x_i$，有
\[
y_i\sim N\bigl(\alpha+\beta(x_i-\bar x),\sigma^2\bigr),
\qquad i=1,\cdots,n.
\]
似然函数为
\[
L(\alpha,\beta)=\prod_{i=1}^n
\frac{1}{\sqrt{2\pi\sigma^2}}
\exp\left[-\frac{\{y_i-\alpha-\beta(x_i-\bar x)\}^2}{2\sigma^2}\right].
\]
取对数，去掉与 $\alpha,\beta$ 无关的常数，最大化 $\ln L$ 等价于最小化
\[
Q(\alpha,\beta)=\sum_{i=1}^n\left[y_i-\alpha-\beta(x_i-\bar x)\right]^2.
\]
分别对 $\alpha,\beta$ 求导：
\[
\frac{\partial Q}{\partial\alpha}
=-2\sum_{i=1}^n\left[y_i-\alpha-\beta(x_i-\bar x)\right]=0.
\]
由于 $\sum_{i=1}^n(x_i-\bar x)=0$，得到
\[
\hat\alpha=\bar y=\frac1n\sum_{i=1}^n y_i.
\]
再对 $\beta$ 求导：
\[
\frac{\partial Q}{\partial\beta}
=-2\sum_{i=1}^n(x_i-\bar x)
\left[y_i-\alpha-\beta(x_i-\bar x)\right]=0.
\]
代入 $\hat\alpha=\bar y$，得
\[
\hat\beta=\frac{\sum_{i=1}^n(x_i-\bar x)(y_i-\bar y)}
{\sum_{i=1}^n(x_i-\bar x)^2}.
\]
若记
\[
s_x^2=\frac1n\sum_{i=1}^n(x_i-\bar x)^2,
\qquad
r_{xy}=\frac1n\sum_{i=1}^n(x_i-\bar x)(y_i-\bar y),
\]
则
\ansmath{\hat\alpha=\bar y,
\qquad
\hat\beta=\frac{r_{xy}}{s_x^2}}
\end{solution}
\end{question}

\begin{question}
静止的双粲重子的衰变时间 $t$ 服从均值为 $\tau$ 的指数分布，其中 $\tau$ 为其固有寿命。假定某实验观测到一个双粲重子事件，且在其质心系的衰变时间为 $0.3\,\mathrm{ps}$。求 $\tau$ 的 $90\%$ 置信水平下的中心置信区间。
\begin{solution}
只有一个观测量，记
\[
t_{\mathrm{obs}}=0.3\,\mathrm{ps}.
\]
指数分布密度为
\[
g(t;\tau)=\frac1\tau e^{-t/\tau},\qquad t\ge0.
\]
令
\[
U=\frac{T}{\tau}.
\]
则 $U$ 服从均值为 $1$ 的指数分布，分布函数为 $F_U(u)=1-e^{-u}$。中心 $90\%$ 区间两端各留 $5\%$ 尾概率，因此
\[
u_{0.05}=-\ln0.95,
\qquad
u_{0.95}=-\ln0.05.
\]
由
\[
u_{0.05}<\frac{t_{\mathrm{obs}}}{\tau}<u_{0.95}
\]
反解得
\[
\frac{t_{\mathrm{obs}}}{u_{0.95}}<\tau<\frac{t_{\mathrm{obs}}}{u_{0.05}}.
\]
代入 $t_{\mathrm{obs}}=0.3\,\mathrm{ps}$，得到
\[
\tau\in\left[\frac{0.3}{-\ln0.05},\frac{0.3}{-\ln0.95}\right]
=[0.100,5.85]~\mathrm{ps}.
\]
\ansmath{\tau\in[0.100,\,5.85]~\mathrm{ps}\quad(90\%\ \mathrm{CL})}
\end{solution}
\end{question}

\begin{question}
假设两个不同的实验对同一个物理量 $x$ 进行了相互独立的测量，得到的测量结果分别为
\[
x_1\pm\sigma_1,
\qquad
x_2\pm\sigma_2.
\]
这两个测量结果完全不相关。用最小二乘法求最佳平均值及其不确定度。
\begin{solution}
两次测量不相关，协方差矩阵为
\[
V=\begin{pmatrix}
\sigma_1^2&0\\
0&\sigma_2^2
\end{pmatrix}.
\]
定义卡方量
\[
\chi^2(x)=\frac{(x_1-x)^2}{\sigma_1^2}
+\frac{(x_2-x)^2}{\sigma_2^2}.
\]
令 $\partial\chi^2/\partial x=0$，得
\[
-2\frac{x_1-x}{\sigma_1^2}
-2\frac{x_2-x}{\sigma_2^2}=0.
\]
因此
\[
x\left(\frac1{\sigma_1^2}+\frac1{\sigma_2^2}\right)
=\frac{x_1}{\sigma_1^2}+\frac{x_2}{\sigma_2^2}.
\]
令
\[
w_i=\frac1{\sigma_i^2},\qquad w=w_1+w_2,
\]
则最佳估计为
\[
\hat x=\frac{w_1x_1+w_2x_2}{w}.
\]
其方差为
\[
\sigma_{\hat x}^2=\frac1w.
\]
所以
\ansmath{\hat x=\frac{x_1/\sigma_1^2+x_2/\sigma_2^2}{1/\sigma_1^2+1/\sigma_2^2},
\qquad
\sigma_{\hat x}=\frac1{\sqrt{1/\sigma_1^2+1/\sigma_2^2}}}
\end{solution}
\end{question}

\begin{question}
假设两个不同的实验对同一个物理量 $x$ 进行了相互独立的测量，得到的测量结果分别为
\[
x_1\pm\sigma_{s1}\pm\sigma_{c1},
\qquad
x_2\pm\sigma_{s2}\pm\sigma_{c2}.
\]
其中第一项不确定度为统计不确定度，两个实验完全不相关；第二项不确定度为系统不确定度，是由相同的外部输入单因素引入的，完全正相关。其他来源引入的系统不确定度可以忽略。求这两个测量结果的最佳平均值及其不确定度。
\begin{solution}
统计不确定度彼此不相关；共同来源的系统不确定度完全正相关。因此协方差矩阵为
\[
V=\begin{pmatrix}
\sigma_{s1}^2+\sigma_{c1}^2 & \sigma_{c1}\sigma_{c2}\\
\sigma_{c1}\sigma_{c2} & \sigma_{s2}^2+\sigma_{c2}^2
\end{pmatrix}.
\]
记
\[
v_1=\sigma_{s1}^2+\sigma_{c1}^2,
\qquad
v_2=\sigma_{s2}^2+\sigma_{c2}^2,
\qquad
c=\sigma_{c1}\sigma_{c2}.
\]
对同一个真值 $x$ 的广义最小二乘估计为
\[
\hat x=\frac{\mathbf 1^{T}V^{-1}\mathbf x}{\mathbf 1^{T}V^{-1}\mathbf 1},
\qquad
\mathbf x=(x_1,x_2)^T,
\qquad
\mathbf 1=(1,1)^T.
\]
对二阶矩阵直接计算可得权重
\[
w_1=\frac{v_2-c}{v_1+v_2-2c},
\qquad
w_2=\frac{v_1-c}{v_1+v_2-2c},
\qquad
w_1+w_2=1.
\]
因此
\[
\hat x=w_1x_1+w_2x_2.
\]
代回 $v_1,v_2,c$，得
\[
\hat x=\frac{(\sigma_{s2}^2+\sigma_{c2}^2-\sigma_{c1}\sigma_{c2})x_1
+(\sigma_{s1}^2+\sigma_{c1}^2-\sigma_{c1}\sigma_{c2})x_2}
{\sigma_{s1}^2+\sigma_{c1}^2+\sigma_{s2}^2+\sigma_{c2}^2-2\sigma_{c1}\sigma_{c2}}.
\]
估计量方差为
\[
\sigma_{\hat x}^2=\frac{1}{\mathbf 1^{T}V^{-1}\mathbf 1}
=\frac{v_1v_2-c^2}{v_1+v_2-2c}.
\]
即
\[
\sigma_{\hat x}^2
=\frac{\sigma_{s1}^2\sigma_{s2}^2+\sigma_{c1}^2\sigma_{s2}^2+\sigma_{c2}^2\sigma_{s1}^2}
{\sigma_{s1}^2+\sigma_{c1}^2+\sigma_{s2}^2+\sigma_{c2}^2-2\sigma_{c1}\sigma_{c2}}.
\]
故
\ansmath{\hat x=\frac{(\sigma_{s2}^2+\sigma_{c2}^2-\sigma_{c1}\sigma_{c2})x_1
+(\sigma_{s1}^2+\sigma_{c1}^2-\sigma_{c1}\sigma_{c2})x_2}
{\sigma_{s1}^2+\sigma_{c1}^2+\sigma_{s2}^2+\sigma_{c2}^2-2\sigma_{c1}\sigma_{c2}},\qquad
\sigma_{\hat x}^2=\frac{\sigma_{s1}^2\sigma_{s2}^2+\sigma_{c1}^2\sigma_{s2}^2+\sigma_{c2}^2\sigma_{s1}^2}
{\sigma_{s1}^2+\sigma_{c1}^2+\sigma_{s2}^2+\sigma_{c2}^2-2\sigma_{c1}\sigma_{c2}}}
\end{solution}
\end{question}

